%% =====================================================================
%% Replicating a paper --- course template (SSRN format)
%% ML & FinTech, National Yang Ming Chiao Tung University
%% Instructor: Huei-Wen Teng
%%
%% Content ported from ../manuscript/main.tex into the course template.
%% Every instruction to you is written in \emph{italics} using \inst{...}.
%% Delete each instruction once you have acted on it.
%% =====================================================================

\documentclass[12pt]{article}

\usepackage{adjustbox,array,amssymb,amsmath,amsfonts,eurosym,geometry,graphicx,
            caption,color,setspace,sectsty,comment,footmisc,natbib,pdflscape,
            subfigure,longtable,lscape,float,multirow,booktabs,makecell,
            hyperref,ulem,appendix}

\normalem
\onehalfspacing

\newtheorem{theorem}{Theorem}
\newtheorem{corollary}[theorem]{Corollary}
\newtheorem{proposition}{Proposition}
\newenvironment{proof}[1][Proof]{\noindent\textbf{#1.} }{\ \rule{0.5em}{0.5em}}

\newtheorem{hyp}{Hypothesis}
\newtheorem{subhyp}{Hypothesis}[hyp]
\renewcommand{\thesubhyp}{\thehyp\alph{subhyp}}

\newcolumntype{L}[1]{>{\raggedright\let\newline\\arraybackslash\hspace{0pt}}m{#1}}
\newcolumntype{C}[1]{>{\centering\let\newline\\arraybackslash\hspace{0pt}}m{#1}}
\newcolumntype{R}[1]{>{\raggedleft\let\newline\\arraybackslash\hspace{0pt}}m{#1}}

\geometry{left=1.0in,right=1.0in,top=1.0in,bottom=1.0in}

%% ---------------------------------------------------------------------
%% Instructions are italic. Delete the \inst{...} once you have done it.
%% ---------------------------------------------------------------------
\newcommand{\inst}[1]{\textit{#1}}

\begin{document}

\begin{titlepage}

\title{Replicating the paper:
       Temporal Relational Ranking for Stock Prediction \citep{feng2019temporal}}

\author{\inst{Your name}\thanks{Department of Information Management and Finance,
        National Yang Ming Chiao Tung University. Email:
        \href{mailto:you@nycu.edu.tw}{you@nycu.edu.tw}}}

\date{\today}
\maketitle

\begin{abstract}
\noindent
\inst{TODO (write last, once the Results section has final numbers): briefly explain what
the paper you are replicating does, what it contributes, and how it bridges the gap left by
the papers before it. Then say why you are interested in this paper. Write it as one
paragraph of continuous prose, not as a list. Aim for 150--200 words. Do not simply copy
the original abstract: the abstract of a replication describes the replication.}

\vspace{0in}
\noindent\textbf{Keywords:} \inst{copy the keywords from the paper you are replicating}

\vspace{0in}
\noindent\textbf{JEL Codes:} \inst{copy the JEL codes from the paper you are
replicating. If the paper does not give any, choose from the JEL classification
system and say in a footnote that you chose them yourself.}

\bigskip
\end{abstract}

\setcounter{page}{0}
\thispagestyle{empty}
\end{titlepage}
\pagebreak

\newpage
\clearpage


\section{Motivation}
\label{sec:motivation}

A stock's future return is usually modeled as a function of its own history: its own past
prices, its own past volumes, its own fundamentals. This is convenient, but it discards
information that is easy to observe and often economically meaningful --- companies in the same
industry, connected by a supply chain, or held together in the same funds tend to move together,
because a shock to one is frequently a shock to the group. A model that treats every stock as an
independent time series is, by construction, blind to that channel.

\citet{feng2019temporal} address this directly. Their Relational Stock Ranking (RSR) model
encodes a graph of relations between stocks --- who shares an industry, who is connected through
Wikidata facts about ownership, supply, and competition --- into a Temporal Graph Convolution
layer, and combines it with a sequential (LSTM-based) encoding of each stock's own price history.
The relational signal is used to re-weight and propagate information between connected stocks at
every training step, and the model is trained to rank stocks by next-day return rather than to
regress the return value directly, since for a long-short or top-$k$ trading strategy the relative
ranking of stocks matters more than the exact predicted return.

This project replicates RSR on the same NASDAQ universe and industry-relation graph used in the
original paper, and then asks a question the original paper does not: what happens to the model's
relational signal itself if the graph is not fixed?

\subsection{Motivation for a dynamic relation graph}

In RSR, the relation graph --- whether industry-based or Wikidata-based --- is constructed once,
from static facts, and held fixed for the entire training and evaluation period. This is a
reasonable simplification: industry classification and ownership facts change slowly. But the
\emph{strength} of a relation between two specific stocks is not static. Two firms in the same
sector can move in lockstep for a stretch driven by a shared macro shock, and largely decouple
once that shock passes; a supply-chain tie can matter enormously around an earnings announcement
and be irrelevant the rest of the year. \citet{cui2023temporal} take one step in this direction by
replacing pairwise relations with group-wise hypergraph structures (industry-belonging,
fund-holding), but the group memberships themselves are still static within their evaluation
window. \citet{jafari2022gcnet} propose a graph convolutional approach to the same prediction
problem but, as of this writing, have not released the training code, so the relation-construction
choices in that line of work are not independently reproducible here. In short, every relational
stock-prediction model we could find fixes \emph{which} stocks are related for the duration of
training; none of them let \emph{how strongly} they are related move over time.

Section~\ref{sec:data} shows this is not a minor omission for this particular dataset: computing
rolling 60-trading-day return correlations across the NASDAQ universe used here, the fraction of
stock pairs whose correlation crosses a fixed threshold swings by roughly six- to seven-fold from
one 60-day window to the next. If that variation carries predictive information beyond what a
static graph already captures, a model that is blind to it is leaving something on the table.

\subsection{Research questions}

\begin{itemize}
  \item \textbf{RQ1 (replication).} Does encoding a relation graph between stocks --- static,
  following \citet{feng2019temporal} --- improve next-day return ranking relative to a baseline
  that treats stocks independently?
  \item \textbf{RQ2 (extension).} Does replacing that static relation graph with one that is
  re-estimated over rolling windows improve ranking performance further, relative to both the
  no-relation baseline and the static-relation model?
\end{itemize}

\subsection{Contribution}

We replicate the RSR baseline and relational models on the official NASDAQ dataset, and implement
a third model, structurally identical to RSR except that its relation graph is fed as a
time-indexed input rather than a training-time constant, so that any performance difference
between it and the static-relation model isolates the effect of allowing the graph to change over
time. At convergence, RQ1 does not hold in the direction the literature would predict --- neither
a static nor a dynamic relation graph lowers mean-squared prediction error relative to the
no-relation baseline in this setting. RQ2's answer is more specific than a yes or no: once a
window-alignment bug in the dynamic graph's construction is fixed (Section~\ref{sec:results}), the
dynamic-relation model outperforms both the baseline and static RSR on simulated trading return,
while still trailing the baseline on top-1 ranking accuracy. We report both the initial reversal
between a short and a converged training run, and this split result across metrics, as the
project's substantive findings, rather than treating either as a failed experiment.

\subsection{Related work}
\label{sec:related}

\textbf{Relation-based stock prediction.} \citet{feng2019temporal} is the paper this project
replicates; Section~\ref{sec:methods} describes its architecture in the detail needed to reproduce
it. \citet{cui2023temporal} extend the same broad idea --- that a stock's neighbors carry
predictive information --- from pairwise relations to group-wise hypergraph relations
(industry-belonging and fund-holding hyperedges), with a hierarchical attention mechanism over
intra-hyperedge, inter-hyperedge, and inter-hypergraph structure. \citet{jafari2022gcnet} instead
use a graph convolutional network directly on a stock-relation graph for price-movement
classification; the authors note in their code release that the implementation was withheld
pending journal review, so we treat this as a design reference rather than a reproducible
baseline. All three approaches share the assumption this project relaxes: that the relation graph,
however it is constructed, does not need to change once training begins.

\textbf{Time-varying information transmission between firms.} A separate empirical-finance
literature documents that the channels these models encode as fixed relations are themselves
time-varying. \citet{evgeniou2023uncovering} show that the cross-section of firm-level return
predictability is both sparse and heterogeneous --- which predictors matter, and how much, differs
across firms and shifts over time --- using a machine-learning framework directly comparable in
spirit to the ranking approach used here. This motivates treating relation strength, like
predictor relevance, as something to be estimated period by period rather than assumed constant.

\inst{TODO: add two more citations from related-literature-shortlist.md (supply-chain /
cross-market lead-lag papers) once their full bibliographic details are confirmed --- the
source list gave titles only, not journal/year, so they are not included here yet to avoid
citing an unverified reference. The template asks for five fundamental papers and five
machine-learning papers; only four are cited so far.}


\section{Data}
\label{sec:data}

We use the NASDAQ portion of the dataset released with \citet{feng2019temporal}
(\texttt{Temporal\_Relational\_Stock\_Ranking}, official repository), which contains: (i) daily
end-of-day features for 1{,}026 NASDAQ tickers --- 5-, 10-, 20-, and 30-day moving averages of
close price plus the raw close-to-close return ratio, already normalized by the stock's maximum
price over the sample --- and (ii) a static industry-relation graph over the same 1{,}026 tickers,
a one-hot encoding of shared-industry membership across 97 industry categories sourced from
sector/industry classifications.

The repository's own documentation describes the underlying raw pull as ``historical (30 years)
End-of-day data \ldots{} of more than 8,000 stocks.'' That describes the unfiltered Google Finance
pull, not the modeling dataset actually used for training: the processed feature files cover
\textbf{1{,}246 trading days per stock, roughly five years}, not thirty, and only 1{,}026 of the
raw universe's stocks survive the paper's own liquidity/history filtering. A reader relying on the
top-level description alone would substantially overestimate both the historical depth and the
breadth of the modeling sample.

\subsection{Dependent and independent variables}

\begin{itemize}
  \item \textbf{Response (continuous).} The next-day return of each stock, which the models
  predict and rank across stocks.
  \item \textbf{Predictors (continuous).} For each stock and day: the 5-, 10-, 20-, and 30-day
  moving averages of the close price and the close-to-close return ratio, each normalized by the
  stock's maximum price over the sample.
  \item \textbf{Relation graph (binary).} A $1{,}026\times1{,}026\times1$ one-hot encoding of
  shared-industry membership across 97 industry categories.
\end{itemize}

\inst{TODO: give each variable an acronym and check that the response definition above matches
the original paper. The dataset is the one released with the original paper, so it is not
expected to differ; say so explicitly if that is right.}

\subsection{Data exploration, cleaning and preprocessing}

\subsubsection{Data quality}

Exploratory analysis (full notebook: \texttt{eda.ipynb}) surfaced two further
gaps between the data and its documentation, beyond the sample-size point above:

\begin{enumerate}
  \item \textbf{Missing values are not \texttt{NaN}.} The preprocessing pipeline
  (\texttt{preprocess/eod.py}) pads missing trading days with a sentinel value of $-1234$ rather
  than a null. Undetected, this sentinel dominates naive summary statistics: raw skewness for
  every feature column comes out around $-23$, which would be an extraordinary departure from
  normality for a bounded price ratio. Once the sentinel is stripped, skewness for every column
  falls to an unremarkable $-0.3$ to $-0.4$. The original $-23$ number says nothing about the
  price data; it is an artifact of an undocumented missing-value code.
  \item \textbf{The relation graph includes non-company tickers.} 156 of the 1{,}026 NASDAQ
  tickers (15.2\%) are recorded with industry \texttt{"n/a"} in the sector/industry relation file.
  Inspecting these tickers shows they are ETFs and index funds (e.g.\ \texttt{QQQ}, \texttt{SHV},
  \texttt{IEF}, \texttt{PFF}) rather than operating companies. An ETF's return is a basket average
  with no industry relation of its own, so its presence in a graph meant to encode
  company-to-company relations is noise rather than signal; 17 of these 156 tickers are fully
  isolated (degree zero) in the industry graph as a direct consequence.
\end{enumerate}

Both issues are corrected identically across every model reported in Section~\ref{sec:methods}:
the $-1234$ sentinel is treated as missing (not as a value) wherever it appears, and results are
reported with the 156 ETF/fund tickers still present in the universe unless noted otherwise, since
removing them changes the size of the relation graph in a way that must be applied consistently
across the static and dynamic conditions --- this is listed as an open item in the project
progress notes (\texttt{PROGRESS.md}).

\subsubsection{A first look at whether the relation graph is really static}

To motivate Section~\ref{sec:methods}'s extension directly from the data, we compute rolling
60-trading-day Pearson correlations of daily returns across the full NASDAQ universe and binarize
each window's correlation matrix at $|\text{corr}| > 0.5$. Across the 21 non-overlapping 60-day
windows spanned by the sample, the resulting edge density --- the fraction of stock pairs counted
as related --- ranges from 3.6\% to 22.5\%, roughly a six-fold swing (Table~\ref{tab:density}).
Whatever a static industry or Wikidata graph is meant to approximate, it is not approximating a
constant.

\begin{table}[h]
\centering
\caption{Edge density of the rolling-correlation relation graph, selected 60-day windows
(NASDAQ, 1{,}026 tickers, $|\text{corr}|>0.5$).}
\label{tab:density}
\begin{tabular}{lccccc}
\toprule
Window & 0 & 4 & 14 & 19 & 20 \\
\midrule
Edge density & 3.6\% & 11.5\% & 22.5\% & 4.6\% & 4.2\% \\
\bottomrule
\end{tabular}
\end{table}


\section{Methodology}
\label{sec:methods}

We train and compare three models, all sharing the same NASDAQ universe, the same 5-, 10-, 20-,
30-day moving-average and close-ratio features, and the same next-day ranking objective. All
three are trained with the same optimizer (Adam), the same train/validation/test split by date
(indices 0--755 / 756--1007 / 1008-- of the 1{,}246 trading days), and the same combined loss:
a mean-squared-error regression term on predicted return plus a pairwise ranking term that
penalizes any pair of stocks whose predicted return order disagrees with their realized order,
weighted by $\alpha$.

\subsection{Benchmark model}

\subsubsection{Baseline: Rank\_LSTM}

The baseline (\texttt{training/rank\_lstm.py} in the official repository) encodes each stock's own
feature history with a single-layer LSTM and predicts next-day return from the final hidden state,
with no information exchanged between stocks. This is the ``stocks are independent'' benchmark
that Section~\ref{sec:motivation} argues is the wrong assumption, and the lower bound RSR is
designed to beat.

\subsubsection{RSR: static relation graph}

The RSR model (\texttt{training/relation\_rank\_lstm.py}) takes the same per-stock representation
--- here, the baseline LSTM's learned sequential embedding, exported after training rather than
downloaded, since the original authors' pretrained embedding file is distributed only via an
external link that we chose not to depend on --- and propagates it across the static industry
relation graph via a learned, softmax-normalized attention weight over each stock's neighbors,
before combining the propagated and original representations into the final prediction. The
relation graph itself, a $1{,}026\times1{,}026\times1$ one-hot encoding, is a training-time
constant: identical on day one of training and day 1{,}246.

\subsection{Experimental Design}

\inst{TODO: give a flow chart of your procedure, from raw data to reported number. A reader
should be able to reproduce your result from that figure alone.}

All three models are trained with $\text{seq}=4$, $\text{unit}=32$, the same Adam optimizer, and
the same combined MSE-plus-ranking loss described above, on the same NASDAQ universe
(1{,}026 tickers, Section~\ref{sec:data}).

\subsubsection{Extension: dynamic relation graph}

Our extension (\texttt{training/dynamic\_relation\_rank\_lstm.py}) is architecturally identical to
RSR --- same attention mechanism, same loss, same embedding input --- with exactly one change: the
relation graph is a placeholder fed at each training step rather than a constant baked into the
computation graph. Which graph is fed depends on which 60-trading-day window the current date
falls in, using the rolling-correlation graphs described in Section~\ref{sec:data} and produced by
\texttt{preprocess/dynamic\_relation.py}. Holding the rest of the architecture fixed is
deliberate: any performance difference between this model and static RSR isolates the effect of
letting the graph vary over time, rather than confounding it with an unrelated architectural
change.

\subsection{Performance measures}

We report three metrics on the held-out test split: \emph{MSE}, the mean-squared error of
predicted return; \emph{mrrt}, the mean reciprocal rank of the true top-return stock (higher
means the model's top pick is more often the stock that actually gained the most); and
\emph{btl}, the simulated return of a strategy that buys the model's top-ranked stock each day
(a simple long-only backtest).


\section{Results}
\label{sec:results}

\inst{TODO: put your numbers and the original paper's numbers side by side in one table. The
tables below report only this replication's own runs.}

\subsection{An early result that did not survive convergence}

Training the static-relation graph model (RSR) to full convergence is substantially more
expensive than the baseline --- each training step processes attention weights over the full
$1{,}026\times1{,}026$ relation matrix, versus a single-stock LSTM forward pass for the baseline
--- so we first compared all three models at a matched, reduced budget of 10 epochs to get a
first read before committing to the full schedule. At 10 epochs, the dynamic-relation model led
both the baseline and static RSR on \emph{both} ranking metrics (mrrt $0.048$ vs.\ $0.038$
baseline and $0.042$ RSR; btl $0.87$ vs.\ $0.67$ baseline and $0.48$ RSR), which read, at the time,
as support for RQ2. We report this result and then show why we do not rely on it.

\subsection{Results at matched convergence (50 epochs)}

\begin{table}[h]
\centering
\caption{Test-set performance, all three models trained to 50 epochs on the same NASDAQ universe
and embedding.}
\label{tab:results50}
\begin{tabular}{lccc}
\toprule
Model & Test MSE & Test mrrt & Test btl \\
\midrule
Baseline (Rank\_LSTM, no relation) & 0.0003774 & \textbf{0.0488} & 1.018 \\
RSR (static industry relation) & 0.0003774 & 0.0276 & \textbf{1.130} \\
Dynamic relation (this project) & 0.0003778 & 0.0288 & 0.861 \\
\bottomrule
\end{tabular}
\end{table}

Trained to the same 50-epoch schedule as the baseline, the picture in
Table~\ref{tab:results50} is materially different from the 10-epoch read above, in two ways.

First, \textbf{MSE is essentially identical across all three models} --- the baseline and static
RSR agree to seven decimal places, and the dynamic model is 0.1\% higher. Neither a static nor a
dynamic relation graph lowers mean-squared prediction error relative to treating stocks
independently once training is run to convergence: RQ1, in the specific sense of ``does a relation
graph reduce MSE,'' is not supported here.

Second, and more strikingly, \textbf{the dynamic-relation model's apparent ranking advantage
reverses}. At 50 epochs, the no-relation baseline has the highest mrrt of the three models by a
wide margin (0.049 vs.\ 0.028--0.029), and static RSR has the highest simulated trading return
(1.13 vs.\ 0.86 for the dynamic model). The dynamic-relation model, the model this project was
built to test, is the \emph{worst} performer on both metrics that matter most for an actual
trading application, once training is run long enough to converge.

\subsection{Why the early signal did not hold}

We considered three candidate explanations for the 10-to-50-epoch reversal:

\begin{enumerate}
  \item \textbf{The rolling-correlation construction is noisy.} A 60-day Pearson correlation
  estimated from daily returns is a high-variance statistic, and binarizing it at a single fixed
  threshold ($|\text{corr}|>0.5$) discards the magnitude information that a softer weighting
  scheme could use.
  \item \textbf{Window boundaries were not aligned with the train/validation/test split.} The
  60-day windows in the result above are computed over the full 1{,}246-day sample without regard
  to the valid\_index~$=756$ / test\_index~$=1008$ boundaries used for evaluation, so a single
  window's relation graph can span dates that fall on both sides of a split.
  \item \textbf{The static graph may already be sufficient at this scale.} With 1{,}026 tickers
  and roughly five years of data, the industry-relation graph may already capture most of the
  exploitable cross-sectional structure, leaving little room for a time-varying signal to add
  value net of the noise it introduces.
\end{enumerate}

We tested explanation (2) directly, since it is the only one fixable without a new modeling
choice: \texttt{preprocess/dynamic\_relation.py} was rewritten to chop the rolling windows
separately within the train, validation, and test date ranges, so no window's relation graph spans
a split boundary, and \texttt{dynamic\_relation\_rank\_lstm.py} was updated to look up the correct
window for a given date from the resulting boundary list rather than assuming windows are spaced
uniformly from day zero. Retraining the dynamic-relation model for 50 epochs with these
split-aligned windows changes the result:

\begin{table}[h]
\centering
\caption{Test-set performance, dynamic-relation model before and after aligning window
boundaries to the train/validation/test split (50 epochs; baseline and RSR rows repeated from
Table~\ref{tab:results50} for comparison).}
\label{tab:results-aligned}
\begin{tabular}{lccc}
\toprule
Model & Test MSE & Test mrrt & Test btl \\
\midrule
Baseline (no relation) & 0.0003774 & \textbf{0.0488} & 1.018 \\
RSR (static relation) & 0.0003774 & 0.0276 & 1.130 \\
Dynamic relation, unaligned windows & 0.0003778 & 0.0288 & 0.861 \\
Dynamic relation, split-aligned windows & 0.0003776 & 0.0311 & \textbf{1.214} \\
\bottomrule
\end{tabular}
\end{table}

Aligning the windows to the evaluation split raises the dynamic model's simulated trading return
from the \emph{lowest} of the three models (0.861) to the \emph{highest} (1.214), and raises mrrt
from 0.0288 to 0.0311 --- still below the baseline's 0.0488, but a real improvement. MSE is
essentially unchanged. Explanation (2) is therefore confirmed as a genuine contributor to the
original result, not merely a plausible-sounding possibility: a design flaw in how the dynamic
graph was constructed, unrelated to whether time-varying relations carry information, was
partially responsible for making the dynamic-relation model look worse than it actually performs.

This leaves a more specific and more interesting pattern than either of our earlier reads: the
split-aligned dynamic-relation model \textbf{beats both the baseline and static RSR on simulated
trading return, but still trails the baseline on top-1 ranking accuracy}. Section~\ref{sec:discussion}
takes this asymmetry as the paper's central empirical finding, rather than treating RQ2 as a single
yes/no question.

We treat the reversal between the 10- and 50-epoch results, and the further change after fixing
the window-alignment bug, as the paper's most important methodological findings, independent of
where RQ2 ultimately lands: an early-training comparison between models of different
computational cost is not a safe substitute for a matched-schedule comparison at convergence, and
a preprocessing detail invisible in the model architecture (how a rolling window is chopped) can
be responsible for a result that otherwise reads as a finding about the model itself.

\subsection{Did the replication succeed?}

Not in the direction \citet{feng2019temporal} report. The original paper's central claim is that
encoding relational information lowers prediction error and improves ranking relative to a
relation-blind baseline; at convergence, our replication finds neither. We see three
non-exclusive reasons this replication may diverge from the original result, none of which we
can rule in or out with a single training run:

\begin{enumerate}
  \item \textbf{Hyperparameters were not tuned to this setting.} We used a single configuration
  ($\text{seq}=4$, $\text{unit}=32$) chosen for consistency between the baseline embedding and the
  relational models, not a value selected by the kind of hyperparameter search the original paper
  likely performed. The original paper's own example commands use larger sequence lengths and
  hidden sizes (e.g.\ $\text{seq}=16$, $\text{unit}=64$ for their published wikidata run).
  \item \textbf{A single training run.} Every number in Table~\ref{tab:results50} comes from one
  run per model with a fixed random seed. Given how much the comparison moved between 10 and 50
  epochs (Section~\ref{sec:results}), a single 50-epoch run is not by itself strong evidence that
  training has fully converged or that the ranking between models is stable; it is one draw, not a
  distribution.
  \item \textbf{We used the pretrained embedding, but exported it ourselves.} Because the
  original authors' pretrained sequential embedding is only available via an external link
  (Section~\ref{sec:methods}), the embedding both RSR and the dynamic model consume here was
  exported from our own baseline run rather than the authors' released artifact, which may not be
  trained under the same conditions as the one used in the original paper's reported numbers.
\end{enumerate}

We do not have evidence to prefer one of these explanations over the others; distinguishing them
is exactly what the robustness checks in the project progress notes (\texttt{PROGRESS.md}; a
second random seed and a brief hyperparameter sweep) are for, and that work is not yet done as of
this draft.

\subsection{Interpreting the dynamic-relation result}
\label{sec:discussion}

RQ2 asked whether letting the relation graph vary over time improves on a static graph. Once the
window/split misalignment identified in Section~\ref{sec:results} is fixed, the answer is no
longer a clean no: the split-aligned dynamic-relation model \textbf{beats both the baseline and
static RSR on simulated trading return (btl: 1.214 vs.\ 1.130 and 1.018), but still trails the
baseline on top-1 ranking accuracy (mrrt: 0.031 vs.\ 0.049)} (Table~\ref{tab:results-aligned}).
This is a more informative result than either of our two earlier reads (favorable at 10 epochs,
unfavorable at 50 epochs before the alignment fix), because it is not just a direction but a
\emph{pattern}: the dynamic graph seems to help the model decide which stocks are worth
overweighting in aggregate --- which is what a long-only backtest on the top-ranked pick rewards
--- without helping it get the single best pick exactly right on any given day, which is what
mrrt strictly measures. A relation graph that shifts over time may be adding information about
which \emph{groups} of stocks are moving together in a given period (useful for a strategy that
buys several top-ranked names) while adding enough noise to individual pairwise rankings to hurt
a metric that only rewards nailing the single top pick.

We do not have direct evidence for this account of the mrrt/btl split beyond the two numbers
themselves, and see it as the most concrete open question this project leaves for future work,
ahead of the two untested candidate explanations from Section~\ref{sec:results}: whether the
rolling-correlation construction is too noisy at the pairwise level even though it is informative
in aggregate (explanation 1), and whether a static industry graph is already close to sufficient
at this scale for the ranking task specifically, even where it is not for the trading-return task
(explanation 3).

It is also worth being explicit about what this result does \emph{not} show: it does not show that
industry relations are the ``correct'' representation and correlation-based dynamic relations are
the ``wrong'' one in some general sense, nor the reverse. It shows that this specific dynamic-graph
construction, once a preprocessing bug is fixed, outperforms a graph that never changes on one
metric and underperforms it on another. A different construction --- correlations estimated with
shrinkage rather than a hard threshold, a graph that updates continuously rather than in 60-day
blocks, or dynamics layered on top of the industry graph rather than replacing it --- could shift
this balance further, and distinguishing ``the idea doesn't help on this metric'' from ``this
implementation doesn't help on this metric'' is future work, not something this draft resolves.

\subsection{Limitations}

Beyond the three points above: results are reported for NASDAQ only, using the industry relation
graph only (the original paper also reports a Wikidata-relation condition, which we did not run);
the 156 ETF/fund tickers identified in Section~\ref{sec:data} remain in the relation graph for
every model reported here, rather than being excluded as their absence of a true industry relation
would suggest; and the dynamic-relation construction was tested with one window length (60 days)
and one threshold ($|\text{corr}|>0.5$), neither of which was tuned.


\section{Conclusion}
\label{sec:conclusion}

This project replicated \citet{feng2019temporal}'s Relational Stock Ranking model on its official
NASDAQ dataset and extended it with a dynamic relation graph, re-estimated from rolling 60-day
return correlations instead of held fixed from static industry facts. At a matched, converged
training budget, neither the static nor the dynamic relation graph improved mean-squared
prediction error relative to a relation-blind baseline. The dynamic-relation extension's picture
on the other two metrics changed twice over the course of this project, and both changes are
findings in their own right: an early, reduced-budget comparison suggested the dynamic model was
best on both ranking accuracy and trading return; a matched 50-epoch comparison reversed that,
making it the worst on both; and fixing a window-alignment bug in how the dynamic graph was
constructed reversed the trading-return result again, leaving the dynamic model \emph{ahead} of
both the baseline and static RSR on simulated trading return while still \emph{behind} the
baseline on top-1 ranking accuracy. Two methodological lessons carry beyond this specific
comparison: first, an early-training read is not a safe substitute for a matched-budget comparison
at convergence when the models being compared have different computational cost; second, a
preprocessing detail with no visible connection to model architecture (how a rolling window is
chopped relative to a data split) can be responsible for a result that otherwise reads as evidence
about the model itself.

Following the course's framing of ``potentials of your projects'' (with modifications: a new
dataset or new methods), the most direct next steps are the remaining robustness checks scoped in
the project progress notes (\texttt{PROGRESS.md}) --- a second random seed, and an alternative
dynamic-relation construction less dependent on a single correlation threshold --- together with
the open question raised in Section~\ref{sec:discussion}: why the split-aligned dynamic graph
helps a portfolio-style metric (btl) while still hurting a single-pick metric (mrrt). Answering
that would say something more specific than ``does a dynamic relation graph help'' --- it would
say what kind of information a time-varying graph adds and what kind of prediction task can use
it. Extending the comparison to the Wikidata relation and to NYSE, both available in the same
official dataset but not run here, is the natural next step after that.


\section*{Declaration}

\subsection*{Declaration of Pedagogical Purpose}

This manuscript is prepared as part of the course project for Machine Learning and
FinTech at National Yang Ming Chiao Tung University, instructed by Professor
Huei-Wen Teng. The work is intended solely for pedagogical purposes, including
learning, discussion, and academic training. It does not represent original research
submitted for journal publication, nor should it be considered professional financial
advice.

\subsection*{Declaration of generative AI and AI-assisted technologies in the writing process}

\inst{Name the tools you actually used. Replace the bracketed placeholders below.}

During the preparation of this work, the authors used [ChatGPT] and [Grok], developed
by [OpenAI] and [xAI] respectively, in order to improve the clarity, precision, and
overall quality of the writing. After using these tools, the authors reviewed and
edited the content as needed and take full responsibility for the content of the
article.


\bibliographystyle{chicago}
\bibliography{ref}


\appendix

\section{\inst{Appendix title}}

\inst{Titles have to be concise. Put here anything a reader needs in order to
reproduce your work but does not need in order to follow the argument. Delete this
appendix if you have nothing to put in it.}


\end{document}
